\noindent\textbf{Notation fixed before any comparison.}
We write $n$ for the manuscript's $\nu$. A type is hyperbolic if $2g-2+n>0$.
All profinite homomorphisms are continuous; quotients of profinite groups are by closed normal subgroups. An outer homomorphism to $H$ is a homomorphism modulo conjugation by $H$, unless a smaller conjugating group is explicitly specified. An identification obtained from a geometric base point is invariant only up to the indicated inner conjugations. We never infer injectivity of a profinite completion map from injectivity of the discrete map.
\begin{center}\small
\begin{tabularx}{\textwidth}{@{}L{.19\textwidth}L{.32\textwidth}Y@{}}
\toprule
Our notation & Manuscript notation and scan & Meaning\\\midrule
$A_{g,n}$, $A^!_{g,n}$ & $A_{g,\nu}$, $A^!_{g,\nu}$; 204--205 & Homeomorphisms preserving the puncture set, respectively each puncture.\\
$M^{\pm}_{g,n}$, $M^+_{g,n}$ & $\Gamma_{g,\nu}$, $\Gamma^+_{g,\nu}$, later $\mathcal T$, $\mathcal T^+$; 205, 226, 328 & Extended mapping class group; orientation-preserving subgroup; punctures may be permuted.\\
$P_{g,n}$ & $\Gamma^{!+}_{g,\nu}$; 205, 226 & Pure orientation-preserving mapping class group. Part F uses $G_n=P_{0,n}$.\\
$\Pi_{g,n}$, $\pi_{g,n}$ & $\Pi_{g,\nu}$, $\widehat\Pi_{g,\nu}$; 221, 328 & Discrete surface fundamental group and its full profinite completion, with peripheral data.\\
$O_{g,n}$ & $\widehat{\widehat{\mathcal T}}$; 321--324 & $\Out_{\mathrm{lac}}(\pi_{g,n})$, not a second application of completion.\\
$\Sigma_{g,n}$ & $\widehat{\mathcal T}^{+}=\Sigma$; 323--324 & Geometric image in $O_{g,n}$; distinguished from the abstract completion $\widehat{M^+_{g,n}}$.\\
$N^{\mathrm{norm}}$, $\mathfrak G^{\mathrm{norm}}$ & $\mathcal N_\Sigma$, $\boldsymbol\Gamma_\Sigma$; 324--326 & $N_{O}(\Sigma)$ and $N_{O}(\Sigma)/\Sigma$, \S26 version.\\
$N^{\mathrm{arith}}$, $\mathfrak G^{\mathrm{arith}}$ & $N_{g,\nu}$, $\Gamma_{g,\nu}$; 370--371 & Arithmetic monodromy image and its quotient by $\Sigma$, \S29 version.\\
$\mathscr M_{g,[n]}$ & $M_{g,\nu}$; 326--329 & Moduli stack of curves with an unordered reduced degree-$n$ puncture divisor. $\mathscr M_{g,n}$ has labelled punctures.\\
$C_n$, $T_n=C_n^{\op}$ & No manuscript symbol; STRATA F.7 & Exit category, respectively degeneration category. Completion retains these variances.\\
$G_K$, $\chi$ & $\mathbb I\Gamma_K$, $\chi$; 327, 323 & $\Gal(\overline K/K)$ and the cyclotomic multiplier.\\\bottomrule
\end{tabularx}
\end{center}
\noindent\small
Later sections may use locally declared synonymous letters when quoting a source. The labels \textbf{Known}, \textbf{Proved}, \textbf{Refuted}, and \textbf{Open} classify the exact numbered assertion. ``Open'' also marks an exact assertion left unverified by this audit; we distinguish that from an established open problem in its accompanying text. A conditional implication can be \textbf{Proved} while its antecedent remains an unrefereed claim. ``Verified'' in the bibliography means that the cited source was opened, not that its entire proof was independently certified.\normalsize
